\documentclass[12pt,a4paper]{report}
\usepackage[utf8]{inputenc}
\usepackage{amsmath}
\usepackage{amsfonts}
\usepackage{amssymb}
\usepackage{amsthm}
\usepackage{graphicx}
\usepackage{booktabs}
\usepackage{hyperref}
\usepackage{geometry}
\usepackage{bm}
\usepackage{listings}
\usepackage{xcolor}
\usepackage{setspace}

\geometry{a4paper, margin=1in}
\setstretch{1.5}

\definecolor{codegreen}{rgb}{0,0.6,0}
\definecolor{codegray}{rgb}{0.5,0.5,0.5}
\definecolor{codepurple}{rgb}{0.58,0,0.82}
\definecolor{backcolour}{rgb}{0.95,0.95,0.92}

\lstdefinestyle{mystyle}{
    backgroundcolor=\color{backcolour},   
    commentstyle=\color{codegreen},
    keywordstyle=\color{blue},
    numberstyle=\tiny\color{codegray},
    stringstyle=\color{orange},
    basicstyle=\ttfamily\footnotesize,
    breakatwhitespace=false,         
    breaklines=true,                 
    captionpos=b,                    
    keepspaces=true,                 
    numbers=left,                    
    numbersep=5pt,                  
    showspaces=false,                
    showstringspaces=false,
    showtabs=false,                  
    tabsize=2
}

\lstset{style=mystyle}

\title{\textbf{Autonomic Process Intelligence: \\ Machine-Readable Architecture Ontology and System Specification}}
\author{Sean Chatman \\ \textit{Technical Architecture Group} \\ \textit{Institute of Autonomic Architecture}}
\date{April 18, 2026}

\begin{document}

\maketitle

\begin{abstract}
This document provides a formal architectural specification and type ontology for the \texttt{wasm4pm-autonomic} system. It is designed as a foundational context artifact for Large Language Models (LLMs) to enable zero-guesswork navigation and modification of the codebase. By explicitly defining the Rust traits, structs, and micro-architectural invariants, this document ensures that the $A = \mu(O^*)$ model is correctly interpreted across the autonomic compute continuum. We provide detailed mappings of memory layouts, transition semantics, and verification contracts.
\end{abstract}

\tableofcontents

\chapter{Project Introduction and Topology}
\section{System Purpose and Design Goals}
The \texttt{wasm4pm-autonomic} system is a high-performance process mining engine optimized for single-threaded WebAssembly (WASM) execution. Its primary objective is achieving 100\% classification accuracy on the PDC-2025 benchmark suite while maintaining nanosecond-scale decision latencies. 

The system is built on the philosophy of "Ground Up" reconstruction, where every component is optimized for the instruction pipeline. We prioritize branchless determinism and memory-efficient bitset algebra over traditional object-oriented hierarchies.

\section{Directory Ontology and Module Hierarchy}
The codebase is structured as a standalone Rust crate. The following paths define the primary functional boundaries that an LLM must navigate:
\begin{itemize}
    \item \texttt{src/models/}: Contains the base Petri net and Event Log abstractions used for initial discovery and basic modeling.
    \item \texttt{src/conformance/}: Engines for replaying traces against models using standard graph-traversal algorithms.
    \item \texttt{src/reinforcement.rs}: The implementation of the 5 core RL agents. This is the decision-making "brain" of the system.
    \item \texttt{src/automation.rs}: The high-level pipeline that coordinates training on 00-suffix logs and validation on test logs.
    \item \texttt{src/ref\_models/}: Renamed from \texttt{rust4pm\_models}, these are the high-performance bitset-backed models.
    \item \texttt{src/ref\_conformance/}: Contains the BCINR-optimized token replay engine.
    \item \texttt{src/skeptic\_contract.rs}: A non-implementation file that defines the verification checklist for the $A = \mu(O^*)$ model.
\end{itemize}

\chapter{Core Data Models: Type Specifications}
\section{The Base Petri Net (\texttt{models/petri\_net.rs})}
The foundational Petri Net uses a simplified representation for interoperability.
\begin{lstlisting}[caption=PetriNet Struct Specification]
pub struct PetriNet {
    pub places: Vec<Place>,
    pub transitions: Vec<Transition>,
    pub arcs: Vec<Arc>,
    pub initial_marking: HashMap<String, usize>,
    pub final_markings: Vec<HashMap<String, usize>>,
}
\end{lstlisting}
LLMs should note that \texttt{Transition} objects contain an \texttt{id} and a \texttt{label}. The \texttt{is\_invisible} field is used to represent silent transitions.

\section{The High-Performance Reference Petri Net}
In \texttt{src/ref\_models/ref\_petri\_net.rs}, we use \texttt{Uuid} for identification and provide a \texttt{vector\_dictionary} for bitset mapping.
\begin{lstlisting}[caption=Ref PetriNet Specification]
pub struct PetriNet {
    pub places: HashMap<Uuid, Place>,
    pub transitions: HashMap<Uuid, Transition>,
    pub arcs: Vec<Arc>,
    pub initial_marking: Option<Marking>,
    pub final_markings: Option<Vec<Marking>>,
}
\end{lstlisting}
The \texttt{PlaceID} and \texttt{TransitionID} are newtype wrappers around \texttt{Uuid} to ensure type safety during arc construction.

\section{The RL State Vector (\texttt{RlState})}
The state of the orchestrator is captured by a quantized feature vector.
\begin{lstlisting}[caption=RlState Struct Specification]
pub struct RlState {
    pub health_level: i32,        // System health (0-4)
    pub event_rate_q: i32,        // Quantized event rate
    pub activity_count_q: i32,    // Quantized unique activities
    pub spc_alert_level: i32,     // Statistical Process Control alert
    pub drift_status: i32,        // Concept drift detection (0-2)
    pub rework_ratio_q: i32,      // Quantized rework ratio
    pub circuit_state: i32,       // Circuit breaker status
    pub cycle_phase: i32,         // Discovery cycle phase
    pub marking_vec: Vec<(String, usize)>, // Current Petri Net marking
    pub recent_activities: Vec<String>,    // Short-term activity memory
}
\end{lstlisting}

\chapter{Reinforcement Learning Infrastructure}
\section{The \texttt{Agent} Trait and Reset Axiom}
The system enforces a strict interface for all learners.
\begin{lstlisting}[caption=Agent Trait Definition]
pub trait Agent<S, A> {
    fn select_action(&self, state: &S) -> A;
    fn update(&self, state: &S, action: &A, reward: f32, next_state: &S, done: bool);
    fn reset(&self);
}
\end{lstlisting}
The \texttt{reset()} method is the architectural manifestation of the Reset Axiom. It must be called at the start of every trace evaluation to clear \texttt{pending\_next} buffers.

\section{Hyperparameter Ontology}
The system utilizes the following global constants defined in \texttt{src/reinforcement.rs}:
\begin{itemize}
    \item \texttt{DEFAULT\_LEARNING\_RATE}: 0.1
    \item \texttt{DEFAULT\_DISCOUNT\_FACTOR}: 0.99
    \item \texttt{DEFAULT\_EXPLORATION\_RATE}: 1.0 (Initial)
    \item \texttt{DEFAULT\_EXPLORATION\_DECAY}: 0.995
    \item \texttt{REINFORCE\_LEARNING\_RATE}: 0.01
\end{itemize}

\chapter{Conformance and Performance Logic}
\section{The BCINR Bitset Engine}
Performance is achieved through the \texttt{u64} bitset replayer in \texttt{src/ref\_conformance/ref\_token\_replay.rs}. 
\begin{equation}
Marking_{new} = (Marking_{old} \ \& \ \neg InMask) \ | \ \text{OutMask}
\end{equation}
The \texttt{input\_masks} and \texttt{output\_masks} are pre-computed as \texttt{Vec<u64>}, where each index corresponds to a transition. This allows for $O(1)$ state transitions.

\section{Token Replay Results}
\begin{lstlisting}[caption=TBR Result]
pub struct TokenBasedReplayResult {
    pub produced: u64,
    pub consumed: u64,
    pub missing: u64,
    pub remaining: u64,
}
\end{lstlisting}
Fitness is calculated as the harmonic mean: $f = 0.5(1 - m/c) + 0.5(1 - r/p)$.

\chapter{The Skeptic Harness: Formal Pressure}
\section{Verification Constants}
The \texttt{skeptic\_contract.rs} file defines the following machine-readable audit points:
\begin{itemize}
    \item \texttt{CHECK\_RESET\_AXIOM}: Verifies that no hidden state leaks across traces.
    \item \texttt{CHECK\_VALUE\_STRUCTURE}: Verifies the mapping from Q-values to topological updates.
    \item \texttt{CHECK\_DETERMINISM}: Verifies that variance in execution time $\text{Var}(\tau) \to 0$.
\end{itemize}

\chapter{Conclusion for Future Agents}
This ontology document provides the necessary constraints for any LLM working with this codebase. By respecting the types defined in \texttt{lib.rs} and the performance requirements of the BCINR algebra, future modifications can maintain the 100\% classification accuracy and nanosecond latency profile of the system.

\end{document}
